\chapter{Undoing a Multiplication}\label{ch:factoring}

\begin{etymology}
\textbf{Factor} comes from the Latin \emph{factor}, ``the one who makes'', from
\emph{facere}, ``to make''. To factor a number is to find the things that make
it: to write it back as the product that produced it.

\textbf{Modulus} is Latin for ``a small measure''. Gauss brought it into
arithmetic in 1801, in the \emph{Disquisitiones Arithmeticae}.
\end{etymology}

There are several problems a quantum computer might solve better than a classical
one, and they are not equally useful. Feynman's original one, simulating molecules
and materials, is the one most likely to pay off. Optimization problems, like
choosing the best investment portfolio, are the ones most often promised and the
least supported by evidence. Factoring is neither. It is the problem that made the
field famous, because it threatens the cryptography the internet runs on, and it is
the one whose cost is best understood. That makes it the best place to start
measuring how far away a useful machine is.

The number 42 can be written as $2 \times 3 \times 7$. Finding those three
numbers is what \emph{factoring} means: given a number, find the numbers that
multiply to give it.

For a computer, the two directions are very different. Multiplying two primes
is trivial: $37 \times 53 = 1961$, and a machine does that with numbers
thousands of digits long in a fraction of a second. Going back, from $1961$ to
$37$ and $53$, is hard. For small numbers you can try divisors one by one. For
numbers with hundreds of digits, nobody knows a fast way.

This chapter is about why that matters, and about the move that made factoring
the most famous problem in quantum computing: Peter Shor showed that factoring
is the same as finding a period.

\section{A lock that anyone can close}

\begin{history}{1976--1977 \textperiodcentered{} Stanford and MIT}
In 1976, Whitfield Diffie and Martin Hellman published ``New Directions in
Cryptography'', proposing that two people could agree on a secret over a public
channel. Their scheme rested on the discrete logarithm: computing $g^a \bmod p$
is easy, recovering $a$ from the result is not. A year later, at MIT, Ron
Rivest, Adi Shamir and Leonard Adleman built a system on a different hard
problem, factoring. RSA was published in 1978. In 1997 the British government
revealed that Clifford Cocks, at GCHQ, had found the same idea in 1973 and kept
it secret.
\end{history}

\begin{concept}{Public-key cryptography}
A pair of keys: a public one that anyone can use to lock a message, and a
private one that unlocks it. In RSA, the public key contains a number
$N = p \times q$, and the private key depends on knowing $p$ and $q$.

\textbf{What it fixes:} two strangers can talk securely without ever meeting to
exchange a secret.

\textbf{What it costs:} all of the security rests on a hard problem that nobody
has proved is hard. If someone finds a fast way to factor, the lock opens.
\end{concept}

\section{How hard is hard?}

The best classical algorithm known for large numbers is the \emph{general number
field sieve}. Its running time grows like
\[
  \exp\!\Big(c\,(\ln N)^{1/3}(\ln\ln N)^{2/3}\Big),
  \qquad c = (64/9)^{1/3} \approx 1.923 .
\]
Read the formula piece by piece. $N$ is the number to factor, and $\ln N$ is
proportional to its number of digits, the real measure of the size of the
input. The main part of the exponent is $(\ln N)^{1/3}$; the
$(\ln\ln N)^{2/3}$ factor grows very slowly. The result is called
\emph{subexponential}: faster than any polynomial in the number of digits,
slower than a pure exponential.

\begin{history}{1988--1993 \textperiodcentered{} Pollard and the Lenstras}
John Pollard proposed the idea of the number field sieve in 1988, in a letter
to Andrew Odlyzko. Arjen Lenstra, Hendrik Lenstra, Mark Manasse and Pollard
turned it into a working algorithm, and in 1990 they used it to factor the ninth
Fermat number, $2^{512}+1$, which had resisted everything else. The general
version, for any number, was published in 1993. Thirty years later it is still
the best classical method for large numbers.
\end{history}

Plugging numbers in shows the shape of the growth. The table drops the
lower-order terms hidden in the exact formula, so only the comparison between
rows means anything.

\begin{center}
\begin{tabular}{lrr}
\hline
number & digits & $\exp(\dots)$ \\
\hline
$42$ & 2 & $\approx e^{3.6}$ \\
$1961$ & 4 & $\approx e^{6.0}$ \\
RSA-250 & 250 & $\approx 10^{24}$ \\
RSA-2048 & 617 & $\approx 10^{35}$ \\
\hline
\end{tabular}
\end{center}

RSA-250 is not a hypothetical row. In February 2020, a team led by Fabrice
Boudot and Paul Zimmermann factored it with the number field sieve, using about
2700 core-years of computation. It is the largest RSA challenge number
factored so far. RSA-2048, the size used in practice, sits roughly eleven orders
of magnitude further. Nobody is going to get there by waiting for faster
classical machines.

\section{Shor's answer}

In 1994 Peter Shor, at Bell Labs, showed that a quantum computer can factor in
polynomial time. His bound for the circuit is
\[
  O\!\big((\ln N)^2 \cdot \ln\ln N \cdot \ln\ln\ln N\big).
\]
The $(\ln N)^2$ term dominates. It comes from modular exponentiation, which
does $O(\ln N)$ multiplications of numbers with $O(\ln N)$ digits. The
$\ln\ln N$ and $\ln\ln\ln N$ terms come from the fast multiplication algorithm
of Schönhage and Strassen used inside each step, and in practice they are small
constants.

There is no exponential anywhere. For RSA-2048, $(\ln N)^2$ is about
$2\times10^{6}$ and the whole expression lands around $10^{7}$. That number counts
elementary operations of an ideal machine. On real hardware every one of those
operations has to be protected against errors, and that is where the real
distance lies (Parts~\ref{part:hardware} and~\ref{part:qec}). In 2019, Craig Gidney and Martin Ekerå
estimated that factoring RSA-2048 would take 20 million noisy physical qubits
running for eight hours. In 2025, Gidney brought that down to under one million
qubits, running for under a week. The largest gate-based machines of 2025 had
on the order of a thousand physical qubits.

So Shor did not make RSA breakable tomorrow. What he showed is that the
difficulty of factoring is a fact about classical computers, not about the
problem. The rest of this chapter is about how he found that out.

\section{Factoring is finding a period}

Shor did not attack factoring head-on. He used a reduction from number theory,
known since Gary Miller's work on primality in 1976: if you can find the
\emph{period} of a certain function, you can factor.

Take the number $N$ to factor and pick some $x$ that shares no divisor with $N$
other than 1 (that is, $\gcd(x, N) = 1$). Now compute the powers of $x$, but
keep only the remainder of the division by $N$ at each step. That is what
``mod $N$'' means: $a \bmod N$ is the remainder of $a$ divided by $N$, and
$a \equiv b \pmod N$ means $a$ and $b$ leave the same remainder. The function is
\[
  f(a) = x^a \bmod N .
\]
Take $N = 15$ and $x = 2$:

\begin{center}
\begin{tabular}{r|rrrrrrrrrrrr}
$a$ & 0 & 1 & 2 & 3 & 4 & 5 & 6 & 7 & 8 & 9 & 10 & 11 \\
\hline
$2^a$ & 1 & 2 & 4 & 8 & 16 & 32 & 64 & 128 & 256 & 512 & 1024 & 2048 \\
$2^a \bmod 15$ & 1 & 2 & 4 & 8 & 1 & 2 & 4 & 8 & 1 & 2 & 4 & 8 \\
\end{tabular}
\end{center}

The powers grow without end, but the remainders repeat: $1, 2, 4, 8$, and again.

\begin{figure}[htbp]
\centering
\begin{tikzpicture}[x=8mm, y=3.2mm]
  \draw[black!40] (-0.3,0) -- (11.6,0);
  \draw[black!40] (-0.3,0) -- (-0.3,9.5);
  \foreach \a/\v in {0/1,1/2,2/4,3/8,4/1,5/2,6/4,7/8,8/1,9/2,10/4,11/8} {
    \draw[hfConcept, line width=1.2pt] (\a,0) -- (\a,\v);
    \fill[hfConcept] (\a,\v) circle (2.2pt);
    \node[font=\sffamily\scriptsize, below] at (\a,0) {\a};
  }
  \foreach \v in {1,2,4,8} \node[font=\sffamily\scriptsize, left] at (-0.3,\v) {\v};
  \foreach \s in {0,4,8} {
    \draw[hfHistory, decorate, decoration={brace, amplitude=4pt}] (\s-0.1,9.8) -- (\s+3.1,9.8)
      node[midway, above=3pt, font=\sffamily\scriptsize, text=hfHistory] {$r = 4$};
  }
  \node[font=\sffamily\small, below=5mm] at (5.5,0) {$a$};
\end{tikzpicture}
\caption{The function $2^a \bmod 15$. The remainders repeat every four steps:
the period is $r = 4$.}
\end{figure}

The table does not show any factor of 15. It shows something else: the
\emph{period} of the function, the smallest $r$ with $x^r \equiv 1 \pmod N$.
Here $r = 4$. Now comes the algebra that turns the period into factors. The
goal is modest: build a product that $N$ divides.

If $x^r \equiv 1 \pmod N$, then $x^r - 1 \equiv 0 \pmod N$, which means $N$
divides $x^r - 1$. So $x^r - 1$ is a bigger number that contains the factors of
$N$. When $r$ is even, it splits as a difference of squares:
\[
  x^r - 1 = \big(x^{r/2} - 1\big)\big(x^{r/2} + 1\big).
\]
$N$ divides the product, so with some luck each factor of $N$ ends up in a
different parenthesis. The greatest common divisor, which Euclid's algorithm
computes fast, pulls them out. With $x = 2$ and $r = 4$, $x^{r/2} = 4$:
\[
  \gcd(4 - 1, 15) = 3, \qquad \gcd(4 + 1, 15) = 5 .
\]
And $15 = 3 \times 5$.

It does not always go this well. If $r$ comes out odd, the difference of squares
is not available. If $x^{r/2} \equiv -1 \pmod N$, all of $N$ falls into one
parenthesis and the gcd returns $N$ itself, which is useless. Take $x = 14$:
$14^2 = 196 \equiv 1 \pmod{15}$, so $r = 2$, but $14 \equiv -1$, and nothing
comes out. In both cases you pick another $x$ and try again. For a number with
at least two distinct odd prime factors, a random $x$ works at least half the
time, so a few attempts are enough. And if the divisor you find is not prime,
you have still split the problem into smaller numbers, and you keep going.

\begin{concept}{Reduction to period finding}
Factoring $N$ becomes: pick $x$, find the period $r$ of $x^a \bmod N$, compute
$\gcd(x^{r/2} \pm 1, N)$.

\textbf{What it fixes:} everything in this recipe is classically easy except one
step, finding $r$. The hard part of factoring is isolated in a single, precise
question.

\textbf{What it costs:} classically, finding $r$ is as hard as factoring, because
the period can be as large as $N$ itself. The reduction only pays off if
something else can find periods fast.
\end{concept}

Periods are what Fourier analysis was made for, and a quantum computer can
perform a Fourier transform on an exponentially large table of values in one go.
That is Shor's algorithm in one sentence. How the table of $x^a \bmod N$ is
written into qubits, and how the transform reads its period, is the subject of
Part~\ref{part:algorithm}.
